\section{Évaluation finale sur 100 questions}\label{sec:finale}

\textbf{Statut : \acompleter.} Le jeu de cent questions (\texttt{eval/answer\_set.json}) est volontairement vide à ce stade :
ses questions et ses points de référence sont rédigés à la main, après le gel des prompts, pour qu'aucun réglage n'ait pu
s'appuyer sur eux. Les résultats ci-dessous seront remplis automatiquement à partir de \texttt{results/answer\_eval\_*}.

\begin{center}\fbox{\parbox{0.8\linewidth}{\centering Tableau à générer : évaluation finale sur le jeu de 100 questions (\textit{à compléter})}}\end{center}


\figfile{f10_final_eval}{Évaluation finale : part des réponses correctes et des réponses contenant une hallucination, par
système.}{fig:final}
\figfile{f11_latency_calls}{Temps et appels LLM par question, par système.}{fig:latence}

\subsection{Protocole}
Chaque élément du jeu indique la réponse attendue : \emph{répondre} (avec les articles de référence et des points de
référence rédigés par un juriste), \emph{clarifier} ou \emph{refuser}. Pour chaque système, le harnais
\texttt{eval/run\_answer\_eval.py} enregistre la réponse, les preuves, les citations, le temps et le nombre d'appels LLM, puis
calcule sans LLM : présence de l'article de référence parmi les preuves et parmi les citations, identifiants cités absents
des preuves, contrôle numérique contre les PDF, et justesse des décisions de clarifier ou de refuser. Un juge LLM
\emph{distinct} du générateur (variable \texttt{JUDGE\_MODEL}, température 0) indique quels points de référence la réponse
couvre, si elle les contredit, et quelles affirmations ne sont pas appuyées par les articles cités. Les intervalles sont
obtenus par \emph{bootstrap} sur les questions ; deux systèmes sont comparés question par question. Un échantillon de
réponses est exporté pour un contrôle humain du juge (accord et kappa de Cohen, \texttt{eval/audit\_agreement.py}).

\subsection{Commandes}
\begin{lstlisting}
python eval/check_answer_set.py
python eval/run_answer_eval.py --system v2 --dry-run
python eval/run_answer_eval.py --system v2
python eval/run_answer_eval.py --system v1.1
python eval/run_answer_eval.py --system v1 --v1-worktree <dossier>
python eval/run_answer_eval.py --summarize results/answer_eval_answer_set_v1.jsonl results/answer_eval_answer_set_v1.1.jsonl results/answer_eval_answer_set_v2.jsonl
python eval/run_answer_eval.py --system v2 --export-audit 20
python eval/audit_agreement.py results/audit_sample_v2.csv
\end{lstlisting}
Le système v1 tourne dans un \texttt{git worktree} de l'étiquette \texttt{v1-baseline-runnable}, avec son ancien index ; le
harnais affiche les commandes de préparation et ne les exécute pas lui-même.

\subsection{Attentes pré-enregistrées}
Le journal \texttt{agent/LEARNING.md} fixe le but de cette évaluation (mesurer v1, v1.1 et v2 : exactitude, taux de fausses
citations, taux d'affirmations non appuyées, qualité des refus, comparaison des langues) mais ne contient pas encore
d'attente chiffrée pour elle. \averifier{} Une attente chiffrée devrait y être écrite avant de lancer l'évaluation.
